\documentclass[11pt]{article}
\usepackage[utf8]{inputenc}
\usepackage[T1]{fontenc}
\usepackage[a4paper,margin=1in]{geometry}
\usepackage{lmodern}
\usepackage{microtype}
\usepackage{amsmath,amssymb,mathtools}
\usepackage{mathrsfs}
\usepackage{tikz-cd}
\usepackage{enumitem}
\usepackage{longtable,booktabs,array}
\usepackage[hidelinks]{hyperref}
\setlength{\parindent}{0pt}
\setlength{\parskip}{0.55em}
\emergencystretch=3em
\hbadness=10000
\newcommand{\K}{K}
\newcommand{\Gr}{\operatorname{Gr}}
\newcommand{\Fil}{\operatorname{Fil}}
\newcommand{\ch}{\operatorname{ch}}
\newcommand{\Todd}{\operatorname{Todd}}
\newcommand{\cl}{\operatorname{cl}}
\newcommand{\Parf}{\operatorname{Parf}}
\newcommand{\Tor}{\operatorname{Tor}}
\newcommand{\Spec}{\operatorname{Spec}}
\newcommand{\RRR}{\mathrm{RRR}}
\providecommand{\codim}{\operatorname{codim}}
\providecommand{\Inv}{\operatorname{Inv}}
\providecommand{\Pic}{\operatorname{Pic}}
\providecommand{\Z}{\mathbb Z}
\providecommand{\Q}{\mathbb Q}
\providecommand{\Ahat}{\widehat A}
\providecommand{\Ahp}{\widehat A^{+}}
\providecommand{\That}{\widehat T}
\providecommand{\rad}{\operatorname{rad}}
\providecommand{\Gl}{\operatorname{Gl}}
\providecommand{\GL}{\operatorname{GL}}
\providecommand{\ch}{\operatorname{ch}}
\providecommand{\Todd}{\operatorname{Todd}}
\providecommand{\Gr}{\operatorname{Gr}}

% Batch 066 macro additions
\providecommand{\Atil}{\widetilde A}
\providecommand{\Ctil}{\widetilde C}
\providecommand{\K}{\mathrm K}
\providecommand{\cC}{\mathcal C}
\providecommand{\cF}{\mathscr F}
\providecommand{\eps}{\varepsilon}
\providecommand{\ellc}{\ell}
\providecommand{\Tor}{\operatorname{Tor}}
\providecommand{\Coker}{\operatorname{Coker}}
\providecommand{\Ker}{\operatorname{Ker}}
\providecommand{\Supp}{\operatorname{Supp}}
\providecommand{\codim}{\operatorname{codim}}
\providecommand{\Gr}{\operatorname{G}}
\providecommand{\rank}{\operatorname{rank}}
\providecommand{\rang}{\operatorname{rang}}
\providecommand{\cO}{\mathcal O}

\begin{document}

\begin{center}
{\Large SGA 6 -- Exposé I}\par
{\large Pages source 133--153}\par
\end{center}


% Additional notation for Expose I continuation
\providecommand{\Cpx}{\mathrm C}
\providecommand{\Kpx}{\mathrm K}
\providecommand{\Dpx}{\mathrm D}
\providecommand{\Ho}{\mathrm H}
\providecommand{\Ob}{\operatorname{ob}}
\providecommand{\qis}{\xrightarrow{\sim}}
\providecommand{\RHom}{\mathrm R\!\Hom}
\providecommand{\Ext}{\operatorname{Ext}}
\providecommand{\Hom}{\operatorname{Hom}}
\providecommand{\Tr}{\operatorname{Tr}}

\providecommand{\Add}{\operatorname{Add}}
\providecommand{\Cart}{\operatorname{Cart}}
\providecommand{\perfamp}{\operatorname{perf.amp}}
\providecommand{\perfdim}{\operatorname{perf.dim}}
\providecommand{\toramp}{\operatorname{tor.amp}}
\providecommand{\torf}{\operatorname{torf}}
\providecommand{\injf}{\operatorname{injf}}
\providecommand{\rg}{\operatorname{rg}}
\providecommand{\is}{\mathrm{is}}
\providecommand{\Mod}{\operatorname{Mod}}


\section*{Exposé I, suite : amplitude parfaite, tor-dimension finie et rang}

\textbf{Remarque 4.12.2.} Soit \(\mathcal C\) une catégorie abélienne (qu'on peut regarder comme fibrée sur un site ponctuel), et soit \(\mathcal C_0\) une sous-catégorie additive strictement pleine. Supposons \(\mathcal C_0\) stable par noyau d'épimorphisme et quasi-relevable dans \(\mathcal C\). Alors un objet acyclique de \(\Kpx^b(\mathcal C_0)\) n'est pas en général nul; donc il n'y a pas d'espoir que le foncteur
\[
  \Kpx^b(\mathcal C_0)\longrightarrow \Dpx(\mathcal C)
\]
soit fidèle. Cependant, si \(\Kpx^{b,\mathrm{ac}}(\mathcal C_0)\) désigne la sous-catégorie épaisse de \(\Kpx^b(\mathcal C_0)\) formée des objets acycliques dans \(\mathcal C\), le foncteur canonique
\[
  \Kpx^b(\mathcal C_0)/\Kpx^{b,\mathrm{ac}}(\mathcal C_0)
  \longrightarrow \Dpx(\mathcal C)
\]
est pleinement fidèle.

Indications : On sait déjà (2.7) que le foncteur
\[
  \Kpx^-(\mathcal C_0)/\Kpx^{-,\mathrm{ac}}(\mathcal C_0)
  \longrightarrow \Dpx^-(\mathcal C)
\]
est pleinement fidèle. Montrer, par exemple à l'aide du critère ([V] 4.2.b)(ii)), que le foncteur
\[
  \Kpx^b(\mathcal C_0)/\Kpx^{b,\mathrm{ac}}(\mathcal C_0)
  \longrightarrow
  \Kpx^-(\mathcal C_0)/\Kpx^{-,\mathrm{ac}}(\mathcal C_0)
\]
est pleinement fidèle.

\textbf{Lemme 4.13.} Soient \(F\in\Ob\Dpx(\mathcal C_S)\), \(a,b,c\in\mathbb Z\) tels que \(a\leq b\) et \(a\leq c\). Si \(F\) est d'amplitude parfaite contenue dans \([a,c]\) (4.7) et si \(\Ho^i(F)=0\) pour \(i>b\), alors \(F\) est d'amplitude parfaite contenue dans \([a,d]\), où \(d=\inf(b,c)\).

\emph{Preuve.} On peut supposer \(b\leq c\), et \(F\) strictement parfait tel que \(F^i=0\) pour \(i\notin[a,c]\). Alors \(F\) est isomorphe à son tronqué \(F'\) défini par
\[
  F'^i=F^i\quad (i<b),\qquad F'^b=Z^b(F),\qquad F'^i=0\quad (i>b).
\]
Mais \(Z^b(F)\) admet une résolution finie droite par des objets de \(\mathcal C_0\), donc est localement objet de \(\mathcal C_0\), d'où l'assertion.

\textbf{Corollaire 4.14.} Soient \(F\in\Ob\mathcal C_S\) et \(n\in\mathbb N\). Conditions équivalentes :
\begin{enumerate}[label=(\roman*)]
\item \(\perfamp(F)\subset[-n,0]\);
\item \(F\) admet localement une résolution gauche de longueur \(n\) par des objets de \(\mathcal C_0\), c'est-à-dire l'ensemble des \(X\) au-dessus de \(S\) tels qu'il existe une suite exacte
\[
  0\longrightarrow L^{-n}\longrightarrow L^{-n+1}\longrightarrow \cdots
  \longrightarrow L^0\longrightarrow F_X\longrightarrow 0,
\]
avec \(L^i\in\Ob\mathcal C_{0,X}\) pour tout \(i\), est un raffinement de \(S\).
\end{enumerate}

\emph{Preuve.} C'est une conséquence immédiate de (4.7) et (4.13).

\textbf{Définition 4.14.1.} Nous dirons que \(F\) est de dimension parfaite \(\leq n\), et nous écrirons \(\perfdim(F)\leq n\), si \(F\) vérifie les conditions équivalentes de (4.14). Nous appellerons dimension parfaite de \(F\) le plus petit entier \(n\) tel que \(\perfdim(F)\leq n\); s'il n'existe pas un tel entier, nous dirons que \(F\) est de dimension parfaite infinie.

Quand \(S\) est le site ponctuel et que \(\mathcal C_0\) est la sous-catégorie pleine de \(\mathcal C\) formée des objets projectifs, la dimension parfaite n'est autre que la dimension projective usuelle.

\textbf{Proposition 4.15.} Soient \(a,b,n\in\mathbb Z\), avec \(a\leq b\), \(n\geq0\), et soit \(F\in\Ob\Dpx(\mathcal C)\) tel que \(F^i=0\) (resp. \(\Ho^i(F)=0\)) pour \(i\notin[a,b]\). Si, pour tout \(i\in[a,b]\), \(F^i\) (resp. \(\Ho^i(F)\)) est de dimension parfaite \(\leq n+(i-a)\), alors \(F\) est d'amplitude parfaite contenue dans \([a-n,b]\). Par suite, si pour tout \(i\in[a,b]\), \(F^i\) (resp. \(\Ho^i(F)\)) est parfait, \(F\) l'est aussi.

\emph{Preuve.} Procéder par récurrence sur \(b\) (\(a\) et \(n\) fixés), en tronquant \(F\) et appliquant (4.11), cf. preuve de (2.11).

\textbf{Lemme 4.16.} Soient \(a,b\in\mathbb Z\), \(a\leq b\), et \(F\in\Ob\Dpx(\mathcal C)\) tel que \(F^i=0\) pour \(i\notin[a,b]\). On suppose que \(\perfamp(F)\subset[a,b]\), et que \(F^i\in\Ob\mathcal C_0\) pour \(i>a\). Alors \(F^a\) est localement objet de \(\mathcal C_0\).

\emph{Preuve.} On a une suite exacte, obtenue en tronquant \(F\),
\[
  0\longrightarrow F'\longrightarrow F\longrightarrow F^a[-a]\longrightarrow 0,
\]
avec \(\perfamp(F')\subset[a+1,b+1]\). De (4.11) et (4.13) on déduit que \(\perfdim(F^a)=0\), c'est-à-dire que \(F^a\) est localement objet de \(\mathcal C_0\).

\textbf{Proposition 4.17.} Soient \(F',F''\in\Ob\Dpx(\mathcal C_S)\), et \(a,b\in\mathbb Z\), avec \(a\leq b\). Considérons les conditions :
\begin{enumerate}[label=(\roman*)]
\item \(\perfamp(F')\subset[a,b]\) et \(\perfamp(F'')\subset[a,b]\);
\item \(\perfamp(F'\oplus F'')\subset[a,b]\).
\end{enumerate}
On a (i) \(\Rightarrow\) (ii), et, si \(\mathcal C_0\) est localement stable par facteurs directs, c'est-à-dire si la relation \(E'\oplus E''\in\Ob\mathcal C_0\), pour \(E',E''\in\Ob\mathcal C_X\) et \(X\in\Ob S\), implique que localement \(E'\) et \(E''\) sont objets de \(\mathcal C_0\), alors on a (i) \(\Leftrightarrow\) (ii).

\emph{Preuve.} L'implication (i) \(\Rightarrow\) (ii) est un cas particulier de (4.11). Prouvons la réciproque. Posons \(F=F'\oplus F''\). Comme \(F\) est pseudo-cohérent, \(F'\) et \(F''\) le sont aussi (2.12). D'autre part, la relation \(\Ho^i(F)=0\) pour \(i\notin[a,b]\) implique \(\Ho^i(F')=\Ho^i(F'')=0\) pour \(i\notin[a,b]\). D'après (2.10 a), on peut donc supposer, quitte à se localiser, que \(F'^i\) (resp. \(F''^i\)) est dans \(\Ob\mathcal C_0\) pour \(i\geq a+1\), et \(F'^i=F''^i=0\) pour \(i>b\), puis, en tronquant, que \(F'^i=F''^i=0\) pour \(i<a\). Comme on a \(\perfamp(F)\subset[a,b]\), on déduit de (4.16) que \(F^a\) est localement objet de \(\mathcal C_0\). Or \(F^a=F'^a\oplus F''^a\), donc \(F'^a\) et \(F''^a\) sont localement objets de \(\mathcal C_0\), puisque \(\mathcal C_0\) est localement stable par facteurs directs.

\textbf{Proposition 4.18.} Soit \((S',\mathcal C',\mathcal C'_0)\) un trio satisfaisant aux mêmes hypothèses (4.0) que \((S,\mathcal C,\mathcal C_0)\). Soit \(f:S'\to S\) un morphisme de sites, défini par \(f^{-1}:S\to S'\), et soit
\[
  u:\Dpx^*(\mathcal C)\longrightarrow \Dpx^*(\mathcal C')\qquad (*=-,+,b)
\]
un foncteur au-dessus de \(f\) (cf. (2.13)). Si \(u\) transforme les objets de \(\mathcal C_0\) en complexes \(\mathcal C'_0\)-parfaits, alors \(u\) définit un foncteur
\[
  \Dpx^*(\mathcal C)_{\mathcal C_0\text{-}\mathrm{parf}}
  \longrightarrow
  \Dpx^*(\mathcal C')_{\mathcal C'_0\text{-}\mathrm{parf}}.
\]

\emph{Preuve.} Soit \(F\in\Ob\Dpx^*(\mathcal C)_{\mathcal C_0\text{-}\mathrm{parf}}\). Montrons que \(u(F)\in\Ob\Dpx^*(\mathcal C')_{\mathcal C'_0\text{-}\mathrm{parf}}\). La question étant locale, on peut supposer \(F\) strictement parfait. Il n'y a plus qu'à dévisser, tenant compte du fait que \(\Dpx^*(\mathcal C')_{\mathcal C'_0\text{-}\mathrm{parf}}\) est une sous-catégorie triangulée de \(\Dpx^*(\mathcal C')\) (4.10).

\textbf{Remarque 4.18.1.} Supposons qu'il existe \(m,n\in\mathbb Z\), \(m\leq n\), tels que \(F\in\Ob\mathcal C_0\) implique
\[
  \mathcal C'_0\text{-}\perfamp(u(F))\subset[m,n].
\]
Alors, si \(F\in\Ob\Dpx^*(\mathcal C)\), on a
\[
  \mathcal C_0\text{-}\perfamp(F)
  \subset[a,b]
  \quad\Rightarrow\quad
  \mathcal C'_0\text{-}\perfamp(u(F))
  \subset[a+m,b+n].
\]
La preuve est laissée en exercice au lecteur (utiliser (4.11)).

\textbf{Corollaire 4.19.} Sous les hypothèses de (2.16), si \(E'\in\Ob\Dpx^-(A,S')\) est parfait en tant qu'objet de \(\Dpx^-(A,S')\) et \(E\in\Ob\Dpx^-(A,S)\) est parfait, alors
\[
  E'\overset{\mathrm L}{\otimes}_A E
\]
est parfait.

\emph{Preuve.} On peut supposer que \(E'\in\Ob\Dpx^-(A,S'_A)\), où \(S'\) est l'objet final de \(S'\), et l'on applique (4.18) à
\[
  u=E'\overset{\mathrm L}{\otimes}_A:
  \Dpx^-(A,S)\longrightarrow \Dpx^-(A,S').
\]
En particulier :

\textbf{Corollaire 4.19.1.} \textbf{a)} Le foncteur
\[
  \mathrm L f^*:\Dpx^-(A,S)\longrightarrow \Dpx^-(A,S')
\]
induit un foncteur
\[
  \mathrm L f^*:\Dpx^-(A,S)_{\mathrm{parf}}
  \longrightarrow
  \Dpx^-(A,S')_{\mathrm{parf}}.
\]

\textbf{b)} Si \(A\) est commutatif, le foncteur
\[
  \overset{\mathrm L}{\otimes}_A:
  \Dpx^-(S)\times_S\Dpx^-(S)
  \longrightarrow \Dpx^-(S)
\]
induit un foncteur
\[
  \overset{\mathrm L}{\otimes}_A:
  \Dpx^-(S)_{\mathrm{parf}}\times_S\Dpx^-(S)_{\mathrm{parf}}
  \longrightarrow \Dpx^-(S)_{\mathrm{parf}}.
\]

\textbf{Remarque 4.19.2.} On peut préciser (4.19) en termes d'amplitude parfaite :
\[
  \perfamp(E)\subset[a,b],\qquad
  \perfamp(E')\subset[a',b']
  \quad\Rightarrow\quad
  \perfamp(E'\overset{\mathrm L}{\otimes}_A E)\subset[a+a',b+b'].
\]
En particulier :
\[
  \perfamp(E)\subset[a,b]
  \quad\Rightarrow\quad
  \perfamp(\mathrm L f^*(E))\subset[a,b].
\]

\section*{5. Tor-dimension finie et perfection}

\textbf{5.0.} Le lecteur se sera sûrement posé la question : que manque-t-il à un complexe pseudo-cohérent \(E\) pour être parfait ? Suffit-il, par exemple, que \(E\) soit à cohomologie bornée ? La réponse est non. Prenons en effet pour \(\mathcal C\) la catégorie des modules sur un anneau commutatif noethérien \(A\), et pour \(\mathcal C_0\) la sous-catégorie pleine des \(A\)-modules projectifs de type fini. N'importe quel \(A\)-module de type fini \(E\) est pseudo-cohérent et à cohomologie bornée, mais \(E\) n'est parfait que si \(E\) est en outre de dimension projective finie (ou, ce qui revient au même, de tor-dimension finie). Cet exemple suggère que la différence entre un complexe pseudo-cohérent et un complexe parfait réside dans une question de dimension cohomologique. C'est ce que nous allons examiner maintenant dans le cas de la catégorie des modules d'un topos annelé (cf. (3.3)). Nous aurons besoin de quelques sorties préliminaires sur la notion de tor-dimension finie. Nous serons brefs, vu que cette question figure déjà dans la littérature : cf. [H] II 4 et SGA 4 XVII 4.1.9.

\textbf{Proposition 5.1.} Soient \((S,A)\) un topos annelé, \(X\in\Ob S\), \(F\in\Ob\Dpx^-(A,X)\), \(m,n\in\mathbb Z\) tels que \(m\leq n\). Conditions équivalentes :
\begin{enumerate}[label=(\roman*)]
\item \(F\) est isomorphe, dans \(\Dpx^-(A,X)\), à un complexe \(F'\) tel que \(F'^i\) soit plat pour tout \(i\) et nul pour \(i\notin[m,n]\);
\item pour tout \(A_X\)-module à droite \(E\), on a
\[
  \Ho^i(E\overset{\mathrm L}{\otimes}_A F)=0\quad\text{pour }i\notin[m,n].
\]
\end{enumerate}

\emph{Preuve.} L'implication (i) \(\Rightarrow\) (ii) est triviale; l'implication (ii) \(\Rightarrow\) (i) résulte, par troncature, du lemme suivant.

\textbf{Lemme 5.1.1.} (cf. (4.16)). Soit \(F\in\Ob\Dpx^-(A,X)\). On suppose que \(F\) satisfait à l'hypothèse (ii) de (5.1), que \(F^i=0\) pour \(i\notin[m,n]\), et que \(F^i\) est plat pour \(i\ne m\). Alors \(F^m\) est plat.

\emph{Preuve.} On a un triangle distingué
\[
\begin{tikzcd}
F' \arrow[rr] && F \arrow[dl]\\
& F^m[-m] \arrow[ul,"+1"] &
\end{tikzcd}
\]
où \(F'\) est à composantes plates et nulles en dehors de l'intervalle \([m+1,n]\) (supposer \(m<n\)). Appliquer le foncteur \(E\overset{\mathrm L}{\otimes}_A\cdot\), pour \(E\) un \(A_X\)-module à droite arbitraire, et écrire la suite exacte de cohomologie.

\textbf{Définition 5.2.} Un objet \(F\) de \(\Dpx^-(A,S)\) satisfaisant aux conditions équivalentes de (5.1) sera dit d'amplitude plate contenue dans l'intervalle \([m,n]\), et nous écrirons
\[
  \toramp(F)\subset[m,n].
\]
Nous dirons que \(F\) est d'amplitude plate finie, ou de tor-dimension finie, s'il existe des entiers \(m\) et \(n\) tels que la condition précédente soit vérifiée.

\textbf{5.3.} Soit \(X\in\Ob S\). Nous noterons \(\Dpx^-(A,X)_{\torf}\) la sous-catégorie pleine de \(\Dpx^-(A,X)\) formée des objets de tor-dimension finie. D'après le critère (ii) de (5.1), c'est une sous-catégorie triangulée de \(\Dpx^-(A,X)\). Pour \(X\) variable, les catégories \(\Dpx^-(A,X)_{\torf}\) forment une sous-\(S\)-catégorie triangulée (1.1.4) de \(\Dpx^-(A,S)\), que nous noterons \(\Dpx^-(S)_{A\text{-}\torf}\). Le foncteur \(\overset{\mathrm L}{\otimes}_A\) induit un \(S\)-foncteur exact
\[
  \overset{\mathrm L}{\otimes}_A:
  \Dpx^b(S_A)\times_S \Dpx(A/S)_{\torf}
  \longrightarrow \Dpx^b(S_{\mathbb Z}).
\]

\textbf{5.4.} Pour qu'un \(A\)-module à gauche \(F\) soit d'amplitude plate contenue dans \([-n,0]\), \(n\in\mathbb N\), il faut et il suffit que \(F\) soit de tor-dimension \(\leq n\) au sens ordinaire, c'est-à-dire que \(F\) admette une résolution gauche de longueur \(n\) par des \(A\)-modules plats.

\textbf{5.5.} Soient \(X\in\Ob S\) et \(m,n\in\mathbb Z\), \(m\leq n\). Pour que \(F\in\Ob\Dpx^-(A,X)\) soit d'amplitude plate contenue dans \([m,n]\), il faut et il suffit qu'il existe une famille couvrante \((X_i\to X)\) telle que \(F|X_i\) soit d'amplitude plate contenue dans \([m,n]\). En particulier, si \(S/X\) est de type fini (SGA 4 VI 1.7), et si \(F\) est localement de tor-dimension finie, alors \(F\) est de tor-dimension finie.

\textbf{Proposition 5.6.} Sous les hypothèses de (2.16), si \(E'\) est d'amplitude plate contenue dans \([m',n']\) en tant qu'objet de \(\Dpx^-(A,S')\) et si \(E\) est d'amplitude plate contenue dans \([m,n]\), alors
\[
  E'\overset{\mathrm L}{\otimes}_A E
\]
est d'amplitude plate contenue dans \([m+m',n+n']\).

\emph{Preuve.} C'est immédiat : utiliser la forme (5.1 (i)) de la notion de tor-dimension finie et le fait que, si \(E\) est un \(A\)-module à gauche plat, \(f^{-1}(E)\) est un \(f^{-1}(A)\)-module à gauche plat (SGA 4 IV App.).

En particulier :

\textbf{Corollaire 5.6.1.} \textbf{a)} Le foncteur
\[
  \mathrm L f^*:\Dpx^-(A,S)\longrightarrow \Dpx^-(A,S')
\]
induit un foncteur
\[
  \mathrm L f^*:\Dpx(A,S)_{\torf}\longrightarrow \Dpx(A,S')_{\torf}.
\]

\textbf{b)} Si \(A\) est commutatif, le foncteur
\[
  \overset{\mathrm L}{\otimes}_A:
  \Dpx^-(S)\times_S\Dpx^-(S)
  \longrightarrow \Dpx^-(S)
\]
induit un foncteur
\[
  \overset{\mathrm L}{\otimes}_A:
  \Dpx(S)_{\torf}\times_S\Dpx(S)_{\torf}
  \longrightarrow \Dpx(S)_{\torf}.
\]

\textbf{Remarque 5.6.2.} Soient \(E\in\Ob\Dpx^-(A,S)\), où \(S\) est l'objet final de \(S\), et \(m,n\in\mathbb Z\), \(m\leq n\). Si \(E\) est d'amplitude plate contenue dans \([m,n]\), il en est de même de \(E_s\), pour tout point \(s\) de \(S\). La réciproque est vraie si \(S\) a assez de points; utiliser la forme (5.1(ii)) de la notion d'amplitude plate et le fait que \(\overset{\mathrm L}{\otimes}_A\) commute au passage aux fibres.

\textbf{Exercices 5.7.} \textbf{a)} Supposons \(A\) commutatif et \((S',A')\) fidèlement plat sur \((S,A)\) (cf. (2.16.2)). Pour que \(E\in\Ob\Dpx^-(A,S)\) soit d'amplitude plate contenue dans \([m,n]\), il faut et il suffit que \(f^*(E)\) le soit.

\textbf{b)} Soit \(F\in\Ob\Dpx^+(A,S)\). Montrer l'équivalence des conditions :
\begin{enumerate}[label=(\roman*)]
\item \(F\) est isomorphe, dans \(\Dpx^+(A,S)\), à un complexe \(F'\) tel que \(F'^i\) soit injectif pour tout \(i\) et nul pour \(i\notin[m,n]\);
\item pour tout \(A_S\)-module à gauche \(E\), on a \(\Ext^i(E,F)=0\) pour \(i\notin[m,n]\).
\end{enumerate}
Si \(F\) satisfait aux conditions précédentes, on dit que \(F\) est d'amplitude injective contenue dans \([m,n]\). La sous-catégorie pleine de \(\Dpx^+(A,S)\) formée des objets d'amplitude injective finie est une sous-catégorie triangulée, notée \(\Dpx^+(A,S)_{\injf}\). Montrer que, si \(A\) est commutatif, le foncteur
\[
  \mathrm R\Hom:\Dpx(S)^\circ\times\Dpx^+(S)\longrightarrow \Dpx(S)
\]
induit un foncteur
\[
  \mathrm R\Hom:\Dpx(S)^\circ_{\torf}\times\Dpx(S)_{\injf}
  \longrightarrow \Dpx(S)_{\injf}.
\]

Nous pouvons maintenant répondre à la question de (5.0) et énoncer le résultat principal de ce numéro.

\textbf{Théorème 5.8.} Soient \((S,A)\) un topos annelé, \(E\in\Ob\Dpx(A,S)\), \(m,n\in\mathbb Z\), avec \(m\leq n\). Conditions équivalentes :
\begin{enumerate}[label=(\roman*)]
\item \(\perfamp(E)\subset[m,n]\) (4.8);
\item \(E\) est \((m-1)\)-pseudo-cohérent (2.15) et l'on a \(\toramp(E)\subset[m,n]\) (5.2).
\end{enumerate}
Comme « parfait » implique évidemment « pseudo-cohérent », on en déduit :

\textbf{Corollaire 5.8.1.} Pour que \(E\) soit parfait (4.8), il faut et il suffit que \(E\) soit pseudo-cohérent (2.15) et localement de tor-dimension finie (5.2).

Pour la démonstration, nous aurons besoin de quelques lemmes.

\textbf{Lemme 5.8.2.} Soit \(B\) un anneau de \(S\). Si \(E\) est un \(A\)-module à gauche, \(G\) un \(B\)-module à droite, et \(F\) un \(A\)-\(B\)-bimodule (à gauche sur \(A\), à droite sur \(B\)), il existe un homomorphisme canonique fonctoriel
\begin{equation}\tag{5.8.2.1}
  \Hom_B(F,G)\otimes_A E
  \longrightarrow
  \Hom_B(\Hom_A(E,F),G).
\end{equation}
Celui-ci est un isomorphisme pour \(G\) injectif et \(E\) de présentation finie.

\emph{Preuve.} Définissons d'abord la flèche. Par la formule d'adjonction chère à Cartan, il revient au même de définir un homomorphisme de \(B\)-modules :
\begin{equation}\tag{*}
  \Hom_B(F,G)\otimes_A E\otimes_{\mathbb Z_S}\Hom_A(E,F)
  \longrightarrow G.
\end{equation}
Or on a un homomorphisme canonique de \(A\)-\(B\)-bimodules
\begin{equation}\tag{**}
  E\otimes_{\mathbb Z_S}\Hom_A(E,F)
  \longrightarrow F
\end{equation}
provenant par adjonction de la flèche identité de \(\Hom_A(E,F)\). On a de même un homomorphisme canonique de \(B\)-modules
\begin{equation}\tag{***}
  \Hom_B(F,G)\otimes_A F\longrightarrow G.
\end{equation}
On définit alors (*) comme le composé de (**) et (***). Pour voir que la flèche (5.8.2.1) est un isomorphisme quand \(G\) est injectif et \(E\) de présentation finie, il suffit de remarquer que, si \(G\) est injectif, la source et le but de la flèche sont des foncteurs exacts à droite en \(E\) et que la flèche induit l'identité pour \(E=A\).

Voici maintenant un cas particulier de (5.8).

\textbf{Lemme 5.8.3.} Soit \(E\) un \(A\)-module à gauche. Conditions équivalentes :
\begin{enumerate}[label=(\roman*)]
\item \(E\) est localement facteur direct d'un module libre de type fini;
\item \(E\) est plat et de présentation finie.
\end{enumerate}

\emph{Preuve} (d'après Bourbaki, Alg. comm., chap. I, §2, Exerc. 15). Il est clair que (i) implique (ii). Prouvons (ii) \(\Rightarrow\) (i). En vertu de (1.3.1), il s'agit de montrer que le foncteur \(\Hom_A(E,\cdot)\) est exact. Soit
\[
  F\longrightarrow F''\longrightarrow 0
\]
une suite exacte de \(A\)-modules à gauche. On doit montrer que la suite
\[
  \Hom_A(E,F)\longrightarrow \Hom_A(E,F'')\longrightarrow 0
\]
est exacte. Il suffit pour cela de prouver que, pour tout \(\mathbb Z_S\)-module injectif \(I\), la suite déduite par application du foncteur \(\Hom_{\mathbb Z_S}(\cdot,I)\) est exacte. Or on a un diagramme commutatif, où les flèches verticales sont les flèches canoniques de (5.8.2) :
\[
\begin{array}{ccccccccc}
0&\to&\Hom_{\mathbb Z_S}(F'',I)\otimes_A E&\to&\Hom_{\mathbb Z_S}(F,I)\otimes_A E\\
&&\downarrow&&\downarrow\\
0&\to&\Hom_{\mathbb Z_S}(\Hom_A(E,F''),I)&\to&\Hom_{\mathbb Z_S}(\Hom_A(E,F),I).
\end{array}
\]
Les flèches verticales sont des isomorphismes en vertu de (5.8.2). La ligne du haut est exacte parce que \(E\) est plat. Donc la ligne du bas est exacte, cqfd.

\emph{Preuve de (5.8).} L'implication (i) \(\Rightarrow\) (ii) résulte trivialement des définitions. Prouvons (ii) \(\Rightarrow\) (i). La question étant locale, on peut supposer \(E\) strictement \((m-1)\)-pseudo-cohérent. Comme d'autre part \(E\) est acyclique en degré \(<m\), la flèche de troncature
\[
\begin{array}{ccccccc}
\cdots&\to&E^{m-1}&\to&E^m&\to&E^{m+1}\to\cdots\\
&&\downarrow&&\downarrow&&\downarrow\\
\cdots&\to&0&\to&E^m/B^m&\to&E^{m+1}\to\cdots
\end{array}
\]
est un quasi-isomorphisme. Notons \(E'\) le complexe tronqué. Par construction, \(E'\) est d'amplitude plate contenue dans \([m,n]\) et \(E'^i\) est libre de type fini pour \(i>m\). Donc, d'après (5.1.1), \(E'^m\) est plat. On a d'autre part un triangle distingué
\[
\begin{tikzcd}
E'' \arrow[rr] && E' \arrow[dl]\\
& E'^m[-m] \arrow[ul,"+1"] &
\end{tikzcd}
\]
où \(E''\) est strictement parfait. D'après (2.6), \(E'^m[-m]\) est \((m-1)\)-pseudo-cohérent, c'est-à-dire, par (2.9), que \(E'^m\) est de présentation finie. On conclut grâce à (5.8.3) que \(E'^m\) est localement facteur direct d'un module libre de type fini, donc que l'on a \(\perfamp(E)\subset[m,n]\).

\textbf{Remarques 5.8.4.} \textbf{a)} Utilisant (2.5), on retrouve le fait que les complexes parfaits forment une sous-catégorie triangulée de \(\Dpx(A,S)\).

\textbf{b)} Un complexe parfait, même à cohomologie bornée, n'est pas nécessairement de tor-dimension finie. C'est cependant le cas si \(S\) est de type fini (5.5). Notons \(\Dpx(A,S)_{\mathrm{parf}}\) la sous-\(S\)-catégorie pleine de \(\Dpx^b(A,S)\) formée des complexes parfaits et de tor-dimension finie. Il résulte aussitôt de (2.16.1 b) et (5.3) que, si \(A\) est commutatif, le foncteur \(\overset{\mathrm L}{\otimes}_A\) induit un foncteur
\[
  \overset{\mathrm L}{\otimes}_A:
  \Dpx(S)_{\mathrm{parf}}\times_S \Dpx^b(S)_{\mathrm{coh}}
  \longrightarrow \Dpx^b(S)_{\mathrm{coh}}.
\]
Cette formule est le point de départ de la théorie des intersections sur les schémas. On notera que par contre le produit tensoriel dérivé de deux complexes pseudo-cohérents à cohomologie bornée n'est pas en général à cohomologie bornée.

\textbf{Corollaire 5.9.} Supposons que \(S=(\mathrm{Ens})\), donc \(A\) est un anneau ordinaire, et que tout \(A\)-module à gauche soit de tor-dimension finie; c'est le cas par exemple si \(A\) est de dimension cohomologique globale finie, voir (MULT VI 2.6) et (EGA \(0_{\mathrm{IV}}\) 17.2.8). Alors on a
\[
\text{(i)}\quad \Dpx^b(A,S)=\Dpx(A,S)_{\torf},
\qquad
\text{et}
\qquad
\text{(ii)}\quad \Dpx^b(A,S)_{\mathrm{coh}}=\Dpx(A,S)_{\mathrm{parf}}.
\]
Autrement dit, pour que \(E\in\Ob\Dpx(A,S)\) soit de tor-dimension finie (resp. de tor-dimension finie et parfait), il faut et il suffit que \(E\) soit à cohomologie bornée (resp. à cohomologie bornée et pseudo-cohérent).

\emph{Preuve.} D'après (5.8), il suffit de prouver (i). Soit \(E\in\Ob\Dpx^b(A,S)\). On peut supposer, quitte à prolonger \(E\) par \(0\), que \(E\in\Ob\Dpx^b(A,S)\), où \(S\) est l'objet final de \(S\). Procédant par récurrence sur le nombre d'objets de cohomologie non nuls de \(E\), on se ramène, par dévissage, au cas où \(E\) est un \(A\)-module ordinaire, et l'on gagne.

On peut énoncer plus généralement :

\textbf{Corollaire 5.10.} Soit \(X\in\Ob S\). Supposons que \(S/X\) ait assez de points (SGA 4 IV 6.4.1) et que, pour tout point \(s\) de \(S/X\), tout \(A_s\)-module à gauche soit de tor-dimension finie. Alors on a
\[
  \Dpx^b(A,X)_{\mathrm{coh}}=\Dpx^b(A,X)_{\mathrm{parf}}.
\]
Autrement dit, pour que \(E\in\Ob\Dpx^b(A,X)\) soit pseudo-cohérent il faut et il suffit que \(E\) soit parfait.

\emph{Preuve.} Soient \(E\in\Ob\Dpx^b(A,X)_{\mathrm{coh}}\) et \(s\) un point de \(S/X\). Il s'agit de montrer que \(E\) est parfait « au voisinage de \(s\) », c'est-à-dire qu'il existe un objet \(s\)-ponctué \(U\) de \(S\) au-dessus de \(X\) tel que \(E|U\) soit parfait. Or, le foncteur \(s^*\) étant exact, on a, d'après (2.16.1 a), \(E_s\in\Ob\Dpx^b(A_s)_{\mathrm{coh}}\). Donc, en vertu de (5.9), \(E_s\) est parfait. L'assertion sera donc conséquence de (3.7.2) et du lemme suivant :

\textbf{Lemme 5.10.1.} Soit \(s\) un point de \(S\) et soit \(F\in\Ob\Dpx(A_s)_{\mathrm{parf}}\). Il existe un objet \(s\)-ponctué \(U\) de \(S\) et \(F'\in\Ob\Dpx(A_U)_{\mathrm{parf}}\) tel que \(F'_s=F\).

\emph{Preuve.} C'est évident si \(F\) est un \(A_s\)-module libre de type fini. Supposons maintenant que \(F\) soit projectif de type fini, donc image d'un projecteur
\[
  p:L\longrightarrow L,
\]
où \(L\) est libre de type fini. Il existe alors un « voisinage » \(U\) de \(s\), un \(A_U\)-module libre \(L'\) et un homomorphisme \(p':L'\to L'\) tels que \(p'_s=p\). Comme \(p_s'^2=p'_s\), on peut supposer, quitte à rétrécir \(U\), que \(p'^2=p'\). Donc \(F'=\Im(p')\) est facteur direct de \(L'\) et vérifie \(F'_s=F\). On généralise immédiatement au cas où \(F\) est strictement parfait, c'est-à-dire un complexe à degrés bornés composé de modules projectifs de type fini.

\textbf{Exemples 5.11.} Le corollaire (5.10) s'applique en particulier au cas où \((S,A)\) est défini par un espace annelé en anneaux locaux réguliers, par exemple un schéma régulier, ou un espace analytique complexe lisse sur \(\mathbb C\).

\section*{6. Rang d'un complexe parfait}

\textbf{6.0.} Soit \((X,\mathcal O_X)\) un espace annelé en anneaux locaux. Il est bien connu que l'on peut attacher à tout faisceau localement libre de type fini \(L\) sur un ouvert \(U\) de \(X\) une fonction \(\rg_U(L)\) localement constante sur \(U\) et à valeurs entières, qu'on appelle le rang de \(L\), la famille des fonctions \(\rg_U(L)\) étant caractérisée par les propriétés suivantes :
\begin{enumerate}[label=(\roman*)]
\item si \(V\subset U\) est une inclusion d'ouverts, on a \(\rg_V(L|V)=\rg_U(L)|V\);
\item pour toute suite exacte de \(\mathcal O_U\)-modules localement libres de type fini
\[
  0\longrightarrow L'\longrightarrow L\longrightarrow L''\longrightarrow 0,
\]
on a \(\rg_U(L)=\rg_U(L')+\rg_U(L'')\);
\item \(\rg_X(\mathcal O_X)=1\).
\end{enumerate}
On se propose, dans ce numéro, d'étendre la notion de rang aux complexes parfaits.

\textbf{Définition 6.1} (cf. SGA 5 VIII 2). Soient \(\mathcal C\) une catégorie additive (resp. triangulée) et \(F\) un groupe abélien. On dit qu'une fonction
\[
  f:\Ob(\mathcal C)\longrightarrow F
\]
est une fonction additive de \(\mathcal C\) dans \(F\) si, pour tous \(L,L',L''\in\Ob(\mathcal C)\) tels que l'on ait \(L\simeq L'\oplus L''\) (resp. un triangle distingué
\[
\begin{tikzcd}
L' \arrow[rr] && L \arrow[dl]\\
& L'' \arrow[ul] &
\end{tikzcd}
\]
), on a \(f(L)=f(L')+f(L'')\).

Les fonctions additives de \(\mathcal C\) dans \(F\) forment un groupe abélien, noté \(\Add(\mathcal C,F)\).

\textbf{6.2.} Soient \(\mathcal C\) une catégorie triangulée, \(f\) une fonction additive de \(\mathcal C\) dans un groupe abélien \(F\). Alors \(f\) induit une fonction additive sur la catégorie additive \(\mathcal C^{\mathrm{ad}}\) sous-jacente à \(\mathcal C\), d'où une inclusion
\[
  \Add(\mathcal C,F)\subset \Add(\mathcal C^{\mathrm{ad}},F).
\]
En particulier, on a \(f(0)=0\) et \(f(L)=f(M)\) si \(L\simeq M\). On a d'autre part
\[
  f(L[n])=(-1)^n f(L)
\]
pour tout \(L\in\Ob(\mathcal C)\) et \(n\in\mathbb Z\).

\textbf{6.3.} Soit \(\mathcal C\) une catégorie additive (resp. triangulée). Le groupe \(\Add(\mathcal C,F)\) dépend fonctoriellement de \(F\). Le foncteur \(\Add(\mathcal C,\cdot)\) est représentable (cf. SGA 5 VIII 2) par un groupe abélien noté \(k(\mathcal C)\), muni d'une application additive
\[
  \cl_{\mathcal C}:\Ob(\mathcal C)\longrightarrow k(\mathcal C).
\]
Si \(\mathcal C\) est additive, on prend pour \(k(\mathcal C)\) le symétrisé du monoïde des classes d'isomorphie d'objets de \(\mathcal C\); si \(\mathcal C\) est triangulée, on prend pour \(k(\mathcal C)\) le quotient du groupe abélien libre engendré par \(\Ob(\mathcal C)\) par les relations identifiant \(L\) à \(L'+L''\) chaque fois qu'on a un triangle distingué
\[
  L'\longrightarrow L\longrightarrow L''\longrightarrow L'[1].
\]

Le groupe \(k(\mathcal C)\) dépend fonctoriellement de \(\mathcal C\) par rapport aux foncteurs additifs (resp. exacts). Si \(u,u':\mathcal C\to\mathcal D\) sont des foncteurs additifs (resp. exacts) isomorphes, alors on a \(k(u)=k(u')\).

Si \(\mathcal C\) est une catégorie triangulée, on a un épimorphisme canonique
\[
  k(\mathcal C^{\mathrm{ad}})\longrightarrow k(\mathcal C)
\]
(cf. (6.2)).

\textbf{Lemme 6.4.} Soit \(\mathcal C\) une catégorie additive. Le foncteur canonique
\[
  \mathcal C\longrightarrow \Kpx^b(\mathcal C)
\]
induit un isomorphisme, fonctoriel en \(\mathcal C\),
\[
  k(\mathcal C)\xrightarrow{\sim} k(\Kpx^b(\mathcal C)).
\]

\emph{Preuve.} C'est une conséquence immédiate de la formule suivante, qui résulte trivialement de (6.2) :
\begin{equation}\tag{6.4.1}
  f(L)=\sum_i (-1)^i f(L^i)
\end{equation}
pour \(L\in\Ob\Kpx^b(\mathcal C)\) et \(f\) une fonction additive sur \(\Kpx^b(\mathcal C)\).

\emph{Remarque.} Nous n'utiliserons pas ici la notation usuelle \(K(\mathcal C)\), vu que \(K(\mathcal C)\) désigne déjà, lorsque \(\mathcal C\) est additive, la catégorie des complexes de \(\mathcal C\) « à homotopie près ».

\textbf{Définition 6.5.} Soient \(S\) une catégorie, \(\mathcal C\) une \(S\)-catégorie additive (resp. triangulée) (1.1.1), \(F\) un préfaisceau abélien sur \(S\). On dit qu'une famille de fonctions
\[
  f_X:\Ob(\mathcal C_X)\longrightarrow F(X),
  \qquad X\in\Ob S,
\]
indexée par \(X\in\Ob S\) est une \(S\)-fonction additive de \(\mathcal C\) dans \(F\) si les conditions suivantes sont réalisées :
\begin{enumerate}[label=(\roman*)]
\item pour tout \(X\in\Ob S\), \(f_X\) est une fonction additive (6.1);
\item pour toute flèche \(u:X\to Y\) de \(S\), le diagramme suivant est commutatif :
\[
\begin{tikzcd}
\Ob\mathcal C_Y \arrow[r,"f_Y"] \arrow[d,"u^*"'] & F(Y) \arrow[d,"F(u)"]\\
\Ob\mathcal C_X \arrow[r,"f_X"'] & F(X).
\end{tikzcd}
\]
\end{enumerate}
Les \(S\)-fonctions additives de \(\mathcal C\) dans \(F\) forment un groupe abélien, noté \(\Add_S(\mathcal C,F)\).

\textbf{Remarques 6.6.} \textbf{a)} Notons \(\mathcal C^{\is}\) la \(S\)-catégorie fibrée dont la fibre en \(X\in\Ob S\) est la catégorie ayant mêmes objets que \(\mathcal C_X\) mais dont les flèches sont les isomorphismes de \(\mathcal C_X\). Notons d'autre part \(\underline F\) la \(S\)-catégorie fibrée dont la fibre en \(X\in\Ob S\) est la catégorie discrète, c'est-à-dire dont les seules flèches sont les identités, \(F(X)\), le foncteur image inverse par \(u\in\mathrm{Fl}(S)\) étant \(F(u)\). Alors une \(S\)-fonction additive de \(\mathcal C\) dans \(F\) n'est autre qu'un \(S\)-foncteur cartésien de \(\mathcal C^{\is}\) dans \(\underline F\), induisant sur chaque fibre une fonction additive au sens (6.1). Il nous sera commode d'interpréter l'additivité de la façon suivante. Supposons, pour fixer les idées, \(\mathcal C\) triangulée. La catégorie
\[
  \Cart_S(\mathcal C^{\is},\underline F)
\]
(resp. \(\Cart_S(\Tr(\mathcal C)^{\is},\underline F)\) (1.1.2)) des \(S\)-foncteurs cartésiens de \(\mathcal C^{\is}\) (resp. \(\Tr(\mathcal C)^{\is}\)) dans \(\underline F\) est la catégorie discrète associée à un groupe abélien, et l'on a un homomorphisme
\[
  \rho:\Cart_S(\mathcal C^{\is},\underline F)
  \longrightarrow
  \Cart_S(\Tr(\mathcal C)^{\is},\underline F),
\]
défini par
\[
  \rho(f)(T)=f(L)-f(L')-f(L'')
\]
pour \(f\in\Cart_S(\mathcal C^{\is},\underline F)\) et \(T=(L'\to L\to L''\to L'[1])\in\Ob\Tr(\mathcal C)\). Dire que \(f\in\Cart_S(\mathcal C^{\is},\underline F)\) est additif signifie que \(\rho(f)=0\), autrement dit on a une suite exacte, canonique et fonctorielle,
\begin{equation}\tag{6.6.1}
  0\longrightarrow \Add_S(\mathcal C,F)
  \longrightarrow \Cart_S(\mathcal C^{\is},\underline F)
  \xrightarrow{\rho}
  \Cart_S(\Tr(\mathcal C)^{\is},\underline F).
\end{equation}

\textbf{b)} Soit \(u:X\to Y\) une flèche de \(S\). Les foncteurs « image inverse par \(u\) », étant tous isomorphes entre eux, définissent (6.3) un unique homomorphisme
\[
  k(Y)\longrightarrow k(X).
\]
Par suite, \(X\mapsto k(X)\) définit un préfaisceau abélien sur \(S\), noté \(k(\mathcal C)\). Une \(S\)-fonction additive de \(\mathcal C\) dans \(F\) s'interprète alors simplement comme un homomorphisme de préfaisceaux
\[
  k(\mathcal C)\longrightarrow F,
\]
autrement dit on a
\[
  \Add_S(\mathcal C,F)=\Hom(k(\mathcal C),F).
\]
Le préfaisceau \(k(\mathcal C)\) dépend bien entendu fonctoriellement de \(\mathcal C\) vis-à-vis des \(S\)-foncteurs additifs (resp. exacts).

\textbf{Proposition 6.7.} Soient \(S\) un site, \(\mathcal C\) une \(S\)-catégorie abélienne plate, \(\mathcal C_0\) une sous-\(S\)-catégorie additive strictement pleine de \(\mathcal C\) satisfaisant aux hypothèses (4.0). Alors le morphisme de préfaisceaux
\[
  k(\mathcal C_0)\longrightarrow k(\Parf(\mathcal C,\mathcal C_0))
\]
induit par l'inclusion de \(\mathcal C_0\) dans \(\Parf(\mathcal C,\mathcal C_0)\), définit un isomorphisme sur les faisceaux associés. En d'autres termes, pour tout faisceau abélien \(F\) sur \(S\), l'homomorphisme de restriction
\[
  \Add_S(\Parf(\mathcal C,\mathcal C_0),F)
  \longrightarrow
  \Add_S(\mathcal C_0,F)
\]
est un isomorphisme.

\emph{Preuve.} Soit \(F\) un faisceau abélien sur \(S\). Montrons que la flèche
\[
  \Add_S(\Parf(\mathcal C,\mathcal C_0),F)
  \longrightarrow
  \Add_S(\mathcal C_0,F)
\]
est un isomorphisme. Celle-ci se factorise en
\[
\Add_S(\Parf(\mathcal C,\mathcal C_0),F)
 \xrightarrow{a}
\Add_S(\Kpx^b(\mathcal C_0),F)
 \xrightarrow{b}
\Add_S(\mathcal C_0,F).
\]
Or, d'après (6.4), la flèche \(b\) est un isomorphisme. On est donc ramené à prouver que \(a\) est un isomorphisme.

Considérons le diagramme commutatif suivant, où les lignes sont exactes (6.6.1), et les flèches verticales sont les flèches de restriction :
\[
\begin{array}{ccccccccc}
0&\to&\Add_S(\Parf(\mathcal C),F)&\to&\Cart_S(\Parf(\mathcal C)^{\is},\underline F)&\to&\Cart_S(\Tr(\Parf(\mathcal C))^{\is},\underline F)\\
&&\downarrow a&&\downarrow&&\downarrow\\
0&\to&\Add_S(\Kpx^b(\mathcal C_0),F)&\to&\Cart_S(\Kpx^b(\mathcal C_0)^{\is},\underline F)&\to&\Cart_S(\Tr(\Kpx^b(\mathcal C_0))^{\is},\underline F).
\end{array}
\]
Comme \(F\) est un faisceau, \(\underline F\) est un champ (G II 2.2.4), donc, d'après (4.10), les deux flèches verticales de droite sont des isomorphismes. Il en est donc de même de \(a\), cqfd.

\textbf{Remarque 6.7.1} (cf. 4.10.1). On peut, dans (6.7), remplacer \(\Parf(\mathcal C,\mathcal C_0)\) par \(\Dpx^*(\mathcal C)_{\mathcal C_0\text{-}\mathrm{parf}}\) (\(*=-\) ou \(b\)).

\textbf{6.8.} Soient \((S,A)\) un topos annelé, \(\mathcal C\) et \(\mathcal C_0\) les \(S\)-catégories fibrées envisagées en (4.8). On a des homomorphismes de préfaisceaux abéliens
\[
  k(\mathcal C_0)
  \xrightarrow{\alpha}
  k(\Dpx^-(A,S)_{\mathrm{parf}})
  \xrightarrow{\beta}
  k(\Parf(A,S)),
\]
qui, d'après (6.7), induisent des isomorphismes sur les faisceaux associés. L'homomorphisme \(\alpha\) est compatible avec les images inverses, c'est-à-dire que, si
\[
  u:(S',A')\longrightarrow(S,A)
\]
est un morphisme de topos annelés, on a un diagramme commutatif, où les flèches verticales sont induites par le foncteur \(\mathrm L u^*\) (cf. (4.19.1)) :
\[
\begin{tikzcd}
 k(\mathcal C_0) \arrow[r] \arrow[d] & k(\Dpx^-(A,S)_{\mathrm{parf}}) \arrow[d]\\
 k(\mathcal C'_0) \arrow[r] & k(\Dpx^-(A',S')_{\mathrm{parf}}),
\end{tikzcd}
\]
\(\mathcal C'_0\) étant définie de manière analogue à \(\mathcal C_0\). Si \(A\) est commutatif, le produit tensoriel dérivé munit \(k(\mathcal C_0)\) et \(k(\Dpx^-(A,S)_{\mathrm{parf}})\) de structures de préfaisceaux d'anneaux (cf. (4.19.1)), et \(\alpha\) est un homomorphisme de préfaisceaux d'anneaux.

\textbf{6.9.} Dans la situation de (6.8), supposons que \((S,A)\) soit localement annelé (cf. (2.15.1)). Alors (2.15.1) pour tout \(X\in\Ob S\), \(\mathcal C_{0,X}\) est la catégorie des \(A_X\)-modules localement libres de type fini. Notons \(\mathcal C_1\) la \(S\)-catégorie fibrée définie par
\[
  X\longmapsto\text{la catégorie des } A_X\text{-modules libres de type fini}.
\]
Il est clair que \(\mathcal C_1\) vérifie les hypothèses (4.0). Un objet de \(\Dpx(S)\) est parfait si et seulement s'il est parfait relativement à \(\mathcal C_1\). On a des homomorphismes de préfaisceaux abéliens
\[
  k(\mathcal C_1)\longrightarrow k(\mathcal C_0)\longrightarrow
  k(\Dpx^-(S)_{\mathrm{parf}})\longrightarrow k(\Parf(S))
\]
(les deux de gauche étant des homomorphismes de préfaisceaux d'anneaux), qui induisent, d'après (6.7), des isomorphismes sur les faisceaux associés. On va montrer que les faisceaux associés sont en fait canoniquement isomorphes au faisceau constant \(\mathbb Z_S\).

\textbf{Lemme 6.10.} Soient \((S,A)\) un topos localement annelé, \(U\) un objet de \(S\) distinct de l'objet vide, \(p,q\in\mathbb N\). Si l'on a
\[
  A_U^p\simeq A_U^q,
\]
alors \(p=q\).

\emph{Preuve.} Par un argument similaire à celui donné dans la preuve de (2.15.1), on se ramène au cas où \((S,A)\) est le topos localement annelé défini par un préschéma, et dans ce cas l'assertion est triviale : prendre un point dans \(U\).

\textbf{Définition 6.11.} On appelle rang, et l'on note \(\rg\), la \(S\)-fonction additive de \(\mathcal C_1\) dans \(\mathbb Z_S\) définie par
\[
  \rg(L)=p\quad\text{si } L\simeq A_X^p,
  \qquad X\ne\varnothing,
\]
\[
  \rg(L)=0\quad\text{si } L\in\Ob\mathcal C_{1,X},
  \qquad X=\varnothing.
\]
En vertu de (6.7), la fonction rang se prolonge de manière unique en une \(S\)-fonction additive de \(\Parf(S)\) dans \(\mathbb Z_S\), que nous appellerons encore rang et noterons \(\rg\).

\textbf{Proposition 6.12.} \textbf{a)} Pour \(L,M\in\Ob\Dpx^-(X)_{\mathrm{parf}}\), avec \(X\in\Ob S\), on a
\[
  \rg(L\overset{\mathrm L}{\otimes} M)=\rg(L)\,\rg(M).
\]

\textbf{b)} Soit
\[
  u:(S',A')\longrightarrow(S,A)
\]
un morphisme de topos localement annelés. Pour \(M\in\Ob\Dpx^-(S)_{\mathrm{parf}}\), on a
\[
  \rg(\mathrm L u^*(M))=u^*(\rg(M)).
\]
En d'autres termes, l'homomorphisme rang
\[
  \rg:k(\Dpx^-(S)_{\mathrm{parf}})\longrightarrow \mathbb Z_S
\]
est un homomorphisme de préfaisceaux d'anneaux, compatible avec les images inverses.

\end{document}
